\chapter{A Concrete Naming Certificate for $\ReductionFuelBoundaryUp$}
\label{ch:concrete-instances-reductionfuelboundary-namecert}
\origin{ai}

The $\ReductionFuelBoundaryUp$ packet is the finite refusal certificate for the fuel boundary named by \autoref{ch:visions-strong-normalization-vs-turing-completeness}. It sits between the total-host packet of \autoref{ch:concrete-instances-step-indexed-total-host-namecert}, the bounded universal trace of \autoref{ch:concrete-instances-bounded-universal-trace-namecert}, and the typed endpoint surface of \autoref{ch:concrete-instances-typedreductionnormalform-namecert}. Its role is narrower than a global normalization theorem: it records the finite fuel row, the displayed reduction trace, the endpoint row, and the timeout refusal that prevents a missing endpoint from being reread as convergence.

\falsifiablePrediction{Deleting the fuel row makes timeout and normal-endpoint packets indistinguishable to the local classifier; the accepted carrier must reject that packet because the boundary can no longer separate finite endpoint readback from fuel exhaustion.}

\independenceWitness{stepindexedtotalhost, boundeduniversaltrace, typedreductionnormalform, cellulartrustsubstrate, znormal: none is bijective because each sibling omits at least one of the fuel boundary, timeout-refusal ledger, endpoint readback, or total-host substrate comparison rows.}

\begin{definition}[ReductionFuelBoundary carrier]
\label{def:reduction-fuel-boundary-carrier}
A $\ReductionFuelBoundaryUp$ carrier is a finite $\BHist$ packet
$$
Q=(H,F,T,E,U,A,X,C,P,N)
$$
with the following displayed rows:
\begin{enumerate}
\item $H$ is the total-host reduction row, read through the host-facing side of $\StepIndexedTotalHostUp$;
\item $F$ is the finite fuel row naming the boundary at which the reduction trace is inspected;
\item $T$ is the bounded trace row, read through $\BoundedUniversalTraceUp$ and tied to the same fuel row $F$;
\item $E$ is the typed endpoint row, read through $\TypedReductionNormalFormUp$ when the trace reaches a displayed normal endpoint;
\item $U$ is the timeout-refusal row, recording that fuel exhaustion is a refused endpoint rather than a proof of nontermination, halting, or convergence;
\item $A$ is the audit row separating normal endpoint, timeout, missing fuel, detached trace, and host-substrate equality requests;
\item $X$ records componentwise $\hsame$ transports for $H,F,T,E,U,A,C,P,N$;
\item $C$ records displayed $\Cont$ routes from host reduction to fuel, from fuel to bounded trace, from trace to endpoint or timeout refusal, and from each barred semantic request to the audit row;
\item $P$ is the $\Pkg$ provenance over \autoref{ch:visions-strong-normalization-vs-turing-completeness}, $\StepIndexedTotalHostUp$, $\BoundedUniversalTraceUp$, $\TypedReductionNormalFormUp$, $\CellularTrustSubstrateUp$, $\ZNormalUp$, $\BHist$, $\hsame$, $\Cont$, $\Pkg$, and $\NameCert$ rows;
\item $N$ is the local $\NameCert_{\ReductionFuelBoundaryUp}$ row.
\end{enumerate}
The local classifier compares accepted packets only through $H,F,T,E,U,A,X,C,P,N$. It has no coordinate for an unbounded reduction, a global halting predicate, a substrate universality transfer, a quotient endpoint, or host equality with the bounded trace.
\end{definition}
\leandef{BEDC.Derived.ReductionFuelBoundaryUp.ReductionFuelBoundaryUp}

\begin{theorem}[ReductionFuelBoundary NameCert obligations]
\label{thm:reduction-fuel-boundary-namecert-obligations}
For the carrier of \autoref{def:reduction-fuel-boundary-carrier}, the local $\NameCert_{\ReductionFuelBoundaryUp}$ surface consists exactly of carrier admission for
$$
Q=(H,F,T,E,U,A,X,C,P,N),
$$
total-host reduction exposure through $H$, finite fuel exposure through $F$, bounded trace exposure through $T$, typed endpoint readback through $E$, timeout-refusal exactness through $U$, audit separation through $A$, componentwise transport through $X$, displayed continuation replay through $C$, package provenance through $P$, and local naming through $N$.
\end{theorem}
\begin{proof}
Project the carrier of \autoref{def:reduction-fuel-boundary-carrier}. The accepted packet displays exactly $H,F,T,E,U,A,X,C,P,N$, so every public read is either one of those rows or a displayed route replaying them.

The endpoint branch is finite. The host reduction row $H$ is read at fuel $F$, the trace row $T$ is bounded by the same fuel, and the endpoint row $E$ is available only when the typed normal-form boundary is displayed.

The timeout branch is also finite. The refusal row $U$ records fuel exhaustion as a boundary outcome, while the audit row $A$ separates timeout from normal endpoint, missing fuel, detached trace, host-substrate equality, quotient endpoint, and unbounded convergence requests.

Transport, continuation replay, provenance, and local naming are the remaining rows $X,C,P,N$. Since no carrier coordinate names an unbounded reduction or global halting predicate, the listed rows exhaust the local NameCert obligations.
\end{proof}
\leanchecked{BEDC.Derived.ReductionFuelBoundaryUp.ReductionFuelBoundary\_namecert\_obligations}

\begin{theorem}[ReductionFuelBoundary total-host nonescape]
\label{thm:reduction-fuel-boundary-total-host-nonescape}
Every accepted $\ReductionFuelBoundaryUp$ packet exports no total-host reduction read except through the finite route
$$
H,F,T,E,U,A,X,C,P,N .
$$
A downstream consumer may read a bounded trace, a typed endpoint when displayed, or a timeout refusal at the same fuel boundary. It may not convert timeout into nontermination, convert endpoint readback into an unbounded convergence theorem, identify host reduction with substrate universality, or replace the audit row by quotient endpoint equality.
\end{theorem}
\leanchecked{BEDC.Derived.ReductionFuelBoundaryUp.TotalHostNonescape.ReductionFuelBoundaryTotalHostNonescape}
\begin{proof}
Begin with the displayed route row $C$. Its admissible paths pass from the host reduction row $H$ to the fuel row $F$, then to the bounded trace row $T$, and then to either the typed endpoint row $E$ or the timeout-refusal row $U$.

The audit row $A$ keeps those outcomes separate. Endpoint readback requires $E$; timeout requires $U$; missing fuel, detached trace, host-substrate equality, quotient endpoint, and unbounded convergence requests are ledgered as refusals rather than exported facts.

The componentwise transport and provenance rows preserve the same boundary. Moving the packet through $X$, replaying it through $C$, and packaging it through $P,N$ cannot create a new coordinate for halting, nontermination, substrate universality transfer, or endpoint quotienting.

Thus the packet realizes the total-host side of bounded substrate computation as finite fuel, trace, endpoint, timeout-refusal, audit, $\hsame$, $\Cont$, $\Pkg$, and $\NameCert$ rows.
\end{proof}

\begin{theorem}[ReductionFuelBoundary fuel-timeout exactness]
\label{thm:reduction-fuel-boundary-fuel-timeout-exactness}
For every accepted $\ReductionFuelBoundaryUp$ carrier
$$
Q=(H,F,T,E,U,A,X,C,P,N),
$$
the timeout refusal row $U$ is exact at the displayed fuel boundary $F$. A consumer may read timeout only through the common route from $H$ to $F$ to $T$ and then to $U$; it may not read timeout from a detached fuel row, from an absent bounded trace, from endpoint equality, or from a host-level nontermination predicate.
\end{theorem}
\leanchecked{BEDC.Derived.ReductionFuelBoundaryUp.FuelTimeoutExactness.ReductionFuelBoundary\_fuel\_timeout\_exactness}
\begin{proof}
Project the carrier rows of \autoref{def:reduction-fuel-boundary-carrier}. The rows $H,F,T,E,U,A,X,C,P,N$ are the whole accepted packet, so the timeout coordinate is not available independently of the fuel and trace coordinates.

Use the NameCert obligations of \autoref{thm:reduction-fuel-boundary-namecert-obligations}. They expose finite fuel through $F$, bounded trace through $T$, typed endpoint readback through $E$, and timeout-refusal exactness through $U$ for the same carrier.

Finally apply the total-host nonescape theorem \autoref{thm:reduction-fuel-boundary-total-host-nonescape}. Its route admits only $H,F,T,E,U,A,X,C,P,N$, so a timeout read must pass through $F$ and $T$ before reaching $U$, while detached fuel, absent traces, endpoint equality, and host-level nontermination have no exported row.
\end{proof}

\begin{theorem}[ReductionFuelBoundary endpoint-timeout disjointness]
\label{thm:reduction-fuel-boundary-endpoint-timeout-disjointness}
For every accepted $\ReductionFuelBoundaryUp$ carrier
$$
Q=(H,F,T,E,U,A,X,C,P,N),
$$
the typed endpoint row $E$ and timeout-refusal row $U$ are disjoint public reads. A downstream consumer may transport either branch through $X$, replay it through $C$, and package it through $P,N$, but it may not identify endpoint readback with timeout, quotient the endpoint, or turn timeout into convergence or divergence.
\end{theorem}
\leanchecked{BEDC.Derived.ReductionFuelBoundaryUp.EndpointTimeoutDisjointness.ReductionFuelBoundary\_endpoint\_timeout\_disjointness}
\begin{proof}
Project the same carrier rows. The endpoint coordinate is $E$, the timeout coordinate is $U$, and the audit coordinate is $A$; these are displayed as separate rows in \autoref{def:reduction-fuel-boundary-carrier}.

The NameCert obligations theorem exposes typed endpoint readback and timeout-refusal exactness as distinct obligations for the same packet. The audit row separates normal endpoint, timeout, missing fuel, detached trace, and host-substrate equality requests.

Transport and replay do not merge the rows. The rows $X$ and $C$ preserve the branch already selected by $E$ or $U$, while $P$ and $N$ record provenance and local naming without adding quotient endpoint equality, convergence, divergence, or host equality.

The total-host nonescape theorem then bars every remaining route. Since all exported reads stay inside $H,F,T,E,U,A,X,C,P,N$, there is no coordinate that identifies endpoint readback with timeout.
\end{proof}

\begin{theorem}[ReductionFuelBoundary obligation closure package]
\label{thm:reduction-fuel-boundary-obligation-closure-package}
For every accepted $\ReductionFuelBoundaryUp$ carrier
$$
Q=(H,F,T,E,U,A,X,C,P,N),
$$
\autoref{def:reduction-fuel-boundary-carrier}, \autoref{thm:reduction-fuel-boundary-namecert-obligations}, \autoref{thm:reduction-fuel-boundary-total-host-nonescape}, \autoref{thm:reduction-fuel-boundary-fuel-timeout-exactness}, and \autoref{thm:reduction-fuel-boundary-endpoint-timeout-disjointness} enumerate the full obligation surface: carrier admission, total-host reduction exposure, finite fuel exposure, bounded-trace exposure, typed-endpoint readback, timeout-refusal exactness, audit separation, componentwise $\hsame$ transport, displayed $\Cont$ replay, $\Pkg$ provenance, local naming, nontrivial carrier presence, and field faithfulness for the same finite packet.
\end{theorem}
\begin{proof}
The carrier definition gives the packet rows $H,F,T,E,U,A,X,C,P,N$. These rows supply the named carrier, host, fuel, trace, endpoint, timeout, audit, transport, continuation, provenance, and naming coordinates.

The NameCert obligations theorem states that every public obligation is read through those rows. The total-host nonescape theorem prevents an external read of halting, nontermination, unbounded convergence, host-substrate equality, or quotient endpoint equality.

The fuel-timeout exactness theorem binds $U$ to the displayed boundary $F$ and bounded trace $T$. The endpoint-timeout disjointness theorem keeps $E$ and $U$ separate under transport, replay, package provenance, and local naming.

Nontriviality and field faithfulness are therefore obligations on this same packet rather than additional semantic assumptions. A witness must preserve exactly the displayed rows, and any faithful field read lands back in one of $H,F,T,E,U,A,X,C,P,N$.
\end{proof}
\leanchecked{BEDC.Derived.ReductionFuelBoundaryUp.ObligationClosurePackage.ReductionFuelBoundary\_obligation\_closure\_package}

\begin{theorem}[ReductionFuelBoundary scoped substrate route]
\label{thm:reduction-fuel-boundary-scoped-substrate-route}
Every scoped read of an accepted $\ReductionFuelBoundaryUp$ carrier factors through the displayed finite route
$$
H,F,T,E,U,A,X,C,P,N .
$$
The route exposes total-host reduction, finite fuel, bounded trace, typed endpoint, timeout refusal, audit separation, componentwise $\hsame$ transport, displayed $\Cont$ replay, $\Pkg$ provenance, and the local $\NameCert$ row. The Lean TasteGate obligations for $\ReductionFuelBoundaryUp$ preserve these same carrier fields when they witness field faithfulness, nontriviality, and the local \texttt{taste\_gate} package.
\end{theorem}
\begin{proof}
The carrier definition lists exactly the rows $H,F,T,E,U,A,X,C,P,N$. A scoped read therefore starts from the finite packet of \autoref{def:reduction-fuel-boundary-carrier}, not from an ambient substrate equality, an unbounded reduction trace, or an endpoint quotient.

The NameCert obligations of \autoref{thm:reduction-fuel-boundary-namecert-obligations} expose those rows as the public surface. The total-host nonescape theorem \autoref{thm:reduction-fuel-boundary-total-host-nonescape} keeps every host-facing read inside the same route, while \autoref{thm:reduction-fuel-boundary-endpoint-timeout-disjointness} prevents the endpoint and timeout branches from being identified under transport or replay.

It remains to align the paper route with the Lean evidence named below the theorem list. The required $\mathsf{FieldFaithful}\ \ReductionFuelBoundaryUp$, $\mathsf{Nontrivial}\ \ReductionFuelBoundaryUp$, and \texttt{taste\_gate} witnesses are obligations on the same finite carrier fields, so the Lean side preserves the carrier rows rather than adding a new substrate read.
\end{proof}
\leanchecked{BEDC.Derived.ReductionFuelBoundaryUp.ReductionFuelBoundaryScopedSubstrateRoute}

The Lean \texttt{lean4/BEDC/Derived/ReductionFuelBoundaryUp/TasteGate.lean} development must include explicit $\mathsf{Nontrivial}\ \ReductionFuelBoundaryUp$ and $\mathsf{FieldFaithful}\ \ReductionFuelBoundaryUp$ instances. These instances must preserve the total-host reduction row, finite fuel row, bounded trace row, endpoint row, timeout-refusal row, audit row, componentwise transports, displayed continuation routes, package provenance, and local naming row.

\begin{theorem}[ReductionFuelBoundary scoped timeout route]
\label{thm:reduction-fuel-boundary-scoped-timeout-route}
For every accepted $\ReductionFuelBoundaryUp$ carrier
$$
Q=(H,F,T,E,U,A,X,C,P,N),
$$
every scoped timeout read factors through the common route
$$
H,F,T,U,A,X,C,P,N,
$$
while every scoped endpoint read factors through
$$
H,F,T,E,A,X,C,P,N .
$$
The audit row $A$ keeps these branches disjoint and bars nontermination, convergence, quotient endpoint equality, and host-level halting reads.
\end{theorem}
\begin{proof}
Project the carrier of \autoref{def:reduction-fuel-boundary-carrier}. The accepted packet displays one total-host row, one fuel row, one bounded trace row, the endpoint row $E$, the timeout row $U$, and the common audit, transport, replay, provenance, and naming rows $A,X,C,P,N$.

Use the NameCert obligations of \autoref{thm:reduction-fuel-boundary-namecert-obligations}. They expose timeout-refusal exactness and endpoint readback as public reads only inside the displayed packet, so neither branch can start from an external halting predicate, detached trace, quotient endpoint, or host-level convergence claim.

The timeout branch is forced by \autoref{thm:reduction-fuel-boundary-fuel-timeout-exactness}. A timeout read must pass from $H$ through the same fuel and trace rows $F,T$ before reaching $U$, and then through $A,X,C,P,N$ for audit separation, transport, replay, provenance, and local naming.

The endpoint branch is separated by \autoref{thm:reduction-fuel-boundary-endpoint-timeout-disjointness}. Endpoint readback uses $E$ instead of $U$, while the audit row $A$ prevents endpoint and timeout from being identified under transport or replay.

Finally apply \autoref{thm:reduction-fuel-boundary-total-host-nonescape}. Since all scoped reads remain inside $H,F,T,E,U,A,X,C,P,N$, the packet exports no coordinate for nontermination, convergence, quotient endpoint equality, or host-level halting.
\end{proof}
\leanchecked{BEDC.Derived.ReductionFuelBoundaryUp.ReductionFuelBoundaryScopedTimeoutRoute}

\begin{theorem}[Reduction fuel boundary scoped route]
\label{thm:reduction-fuel-boundary-scoped-route}
Every accepted $\ReductionFuelBoundaryUp$ consumer must expose the bounded trace, step-indexed total host, typed normal form, $\ZNormalUp$ boundary, timeout ledger, endpoint-separation audit, and structural $\BHist$, $\hsame$, $\Cont$, $\Pkg$, and $\NameCert$ rows before reading the fuel boundary. The boundary read is therefore a scoped route through the displayed carrier rows of \autoref{def:reduction-fuel-boundary-carrier}, the obligation package of \autoref{thm:reduction-fuel-boundary-obligation-closure-package}, the substrate route of \autoref{thm:reduction-fuel-boundary-scoped-substrate-route}, and the timeout route of \autoref{thm:reduction-fuel-boundary-scoped-timeout-route}; it is not an unbounded halting predicate, a convergence theorem, a quotient endpoint, or a host-substrate equality.
\end{theorem}
\leanchecked{BEDC.Derived.ReductionFuelBoundaryUp.ReductionFuelBoundaryScopedRoute}
\begin{proof}
Project the carrier rows. The total-host row $H$, fuel row $F$, bounded trace row $T$, typed endpoint row $E$, timeout row $U$, audit row $A$, transport row $X$, continuation row $C$, package row $P$, and local naming row $N$ are the complete classifier-visible surface.

The sibling route theorems fix the dependencies. The obligation package enumerates carrier admission and field faithfulness for those rows, while the substrate and timeout route theorems keep the endpoint and timeout branches inside the same finite packet.

The remaining required names are structural rows of that packet. The $\ZNormalUp$ boundary is read only as part of the typed endpoint surface and audit separation, and the $\BHist$, $\hsame$, $\Cont$, $\Pkg$, and $\NameCert$ rows preserve the displayed route rather than adding a semantic shortcut.

Downstream residual-normalization and candidate-mediated SN consumers, including the parallel-diamond frontier of \autoref{ch:concrete-instances-metacic-parallel-diamond-frontier-namecert}, may read the fuel boundary only after the bounded trace, typed endpoint, timeout refusal, and audit rows are displayed. Thus any scoped consumer must read the listed rows before the fuel boundary can be consumed.
\end{proof}

\begin{theorem}[Reduction fuel boundary public certificate]
\label{thm:reduction-fuel-boundary-public-certificate}
The bounded trace row, typed endpoint row, timeout-refusal row, endpoint-separation audit, substrate route, scoped timeout route, and local $\NameCert_{\ReductionFuelBoundaryUp}$ rows form the public certificate exported by $\ReductionFuelBoundaryUp$. Every public consumer of the boundary must read the carrier, obligation package, substrate route, timeout route, and scoped route before using the fuel boundary as a certificate; the export is not an unbounded halting predicate, a global convergence theorem, a quotient endpoint, or a host-substrate equality.
\end{theorem}
\begin{proof}
Start with the carrier and obligation package, \autoref{def:reduction-fuel-boundary-carrier} and \autoref{thm:reduction-fuel-boundary-obligation-closure-package}. They expose the total-host row, finite fuel row, bounded trace row, typed endpoint row, timeout-refusal row, endpoint-separation audit, transport, replay, provenance, and local naming rows as one finite packet.

Read the substrate route next. By \autoref{thm:reduction-fuel-boundary-scoped-substrate-route}, every boundary consumer stays inside the displayed carrier fields and the same field-faithful substrate evidence; no independent host-substrate equality is exported.

Then read the timeout route and scoped route, \autoref{thm:reduction-fuel-boundary-scoped-timeout-route} and \autoref{thm:reduction-fuel-boundary-scoped-route}. They force endpoint and timeout consumers through the same bounded trace and audit surface before the boundary row is public.

Thus the public certificate is exactly the finite route through bounded trace, typed endpoint, timeout refusal, endpoint separation, substrate, scoped timeout, and local NameCert rows. The refused endpoints have no carrier coordinate and cannot be reconstructed from the certificate.
\end{proof}

\closureat{ReductionFuelBoundaryUp}{publicStr}

\begin{closurestatus}{\ReductionFuelBoundaryUp}
  \constructivestory{A $\ReductionFuelBoundaryUp$ packet is generated from a total-host reduction row, finite fuel row, bounded trace row, typed endpoint row, timeout-refusal row, audit row, componentwise hsame transports, displayed Cont routes, package provenance, and a local NameCert row.}
  \theoryclosure{\publicClosure}
  \scopeclosed{\autoref{def:reduction-fuel-boundary-carrier}, \autoref{thm:reduction-fuel-boundary-obligation-closure-package}, \autoref{thm:reduction-fuel-boundary-scoped-substrate-route}, \autoref{thm:reduction-fuel-boundary-scoped-timeout-route}, \autoref{thm:reduction-fuel-boundary-scoped-route}, and \autoref{thm:reduction-fuel-boundary-public-certificate} bind the finite fuel, bounded trace, typed endpoint, timeout refusal, audit, transport, replay, provenance, local naming, and public certificate rows for the boundary packet.}
  \formalstatus{\unformalizedV}
  \bridgestatus{none}
  \origin{ai}
  \notclaimed{This chapter does not prove Lean verification, bridge checking, unbounded halting, global convergence, nontermination, quotient endpoint equality, or transfer of substrate universality to the total host.}
  \upgradepath{Scoped closure requires Lean coverage for the carrier rows, timeout exactness, endpoint-timeout separation, nontriviality, and field faithfulness for the same packet.}
\end{closurestatus}
